\documentclass[11pt]{article}
\usepackage[margin=1in]{geometry}
\usepackage{graphicx}
\usepackage{amsmath}
\usepackage{booktabs}
\usepackage{hyperref}
\usepackage{float}

\title{PETSc-Based Numerical Methods for PDEs:\\
ParaView-Driven Visualization Workflow}
\author{APMA 4302 Project}
\date{\today}

\begin{document}
\maketitle

\begin{abstract}
This report presents a PETSc-first workflow for solving representative PDE classes
(elliptic, parabolic, and advection-diffusion), validating numerical behavior, and
exporting ParaView-ready data for high-quality 3D visualization.
\end{abstract}

\section{Introduction}
This project is organized around a PETSc-first workflow: every numerical run must
produce ParaView-readable outputs before additional analysis is considered complete.
This constraint keeps the computational and visualization pipelines aligned and
ensures that each PDE experiment yields both quantitative metrics and interpretable
field plots.

\section{Mathematical Models}
\subsection{Poisson Equation}
\[
 -\Delta u = f \quad \text{in } \Omega,\qquad u=0 \text{ on } \partial\Omega.
\]
\subsection{Heat Equation}
\[
 u_t - \kappa \Delta u = g.
\]
\subsection{Advection-Diffusion}
\[
 u_t + \mathbf{v}\cdot\nabla u - \nu\Delta u = s.
\]

\section{Numerical Discretization and Solvers}
We use a structured finite-difference grid on the unit square.
The Poisson problem is discretized with the standard 5-point Laplacian.
Heat and advection-diffusion updates are advanced with backward Euler in time,
leading to a linear system at each step solved through PETSc KSP.

Primary solver settings:
\begin{itemize}
  \item Poisson: CG + Jacobi preconditioning.
  \item Heat: CG + Jacobi preconditioning.
  \item Advection-diffusion: default \texttt{cg} + Jacobi in the project drivers for parallel-safe implicit steps (GMRES + Jacobi when nonsymmetry is stronger; avoid parallel ILU on distributed matrices).
\end{itemize}

\section{Implementation Details}
\begin{itemize}
  \item Environment alignment with \texttt{PETSC\_DIR} and \texttt{PETSC\_ARCH=apma4302-pkgs-opt}.
  \item Data layout and MPI distribution.
  \item Export format \texttt{.vtp} (quad surfaces / polylines) and \texttt{.pvd} time collections for ParaView, under \texttt{output/paraview/vtp/}.
\end{itemize}

\section{Verification and Convergence}
For Poisson verification, the numerical solution is compared to the analytic
reference \(u(x,y)=\sin(\pi x)\sin(\pi y)\). The spatial error decreases
consistently under mesh refinement, as shown in Figure~\ref{fig:spatial-conv}.

\begin{figure}[h]
  \centering
  \includegraphics[width=0.72\textwidth]{visuals/spatial_convergence_loglog.png}
  \caption{Poisson spatial convergence on progressively finer grids.}
  \label{fig:spatial-conv}
\end{figure}

\section{Performance Evaluation}
Strong scaling was measured with repeated runs on the same \(1201\times 1201\)
problem to quantify runtime variability.

\begin{table}[h]
  \centering
  \begin{tabular}{rrrrrr}
    \toprule
    Ranks & Count & Mean time (s) & Std (s) & Speedup & Efficiency \\
    \midrule
    1 & 6 & 0.02453 & 0.00188 & 1.00 & 1.00 \\
    2 & 5 & 0.01620 & 0.00237 & 1.51 & 0.76 \\
    4 & 5 & 0.01348 & 0.00448 & 1.82 & 0.45 \\
    8 & 5 & 0.04346 & 0.03091 & 0.56 & 0.07 \\
    \bottomrule
  \end{tabular}
  \caption{Strong-scaling summary from repeated runs (\texttt{strong\_scaling\_summary.csv}).}
\end{table}

\begin{figure}[h]
  \centering
  \includegraphics[width=0.72\textwidth]{visuals/strong_scaling_speedup_errorbars.png}
  \caption{Strong-scaling speedup with mean \(\pm\) one standard deviation.}
\end{figure}

\begin{figure}[h]
  \centering
  \includegraphics[width=0.72\textwidth]{visuals/strong_scaling_runtime_errorbars.png}
  \caption{Runtime variability by MPI rank, showing overhead-dominated behavior at 8 ranks.}
\end{figure}

\section{Visualization Results}
\subsection{Poisson}
\begin{figure}[h]
  \centering
  \includegraphics[width=0.8\textwidth]{visuals/poisson_solution_surface.png}
  \caption{Poisson solution field rendered in ParaView.}
\end{figure}

\subsection{Heat}
\begin{figure}[h]
  \centering
  \includegraphics[width=0.72\textwidth]{visuals/heat_t0.png}
  \caption{Representative heat-equation snapshot exported from ParaView.}
\end{figure}

\subsection{Advection-Diffusion}
\begin{figure}[H]
  \centering
  \includegraphics[width=0.72\textwidth]{visuals/advection_diffusion_volume.png}
  \caption{Advection-diffusion concentration at a representative late-time snapshot.}
\end{figure}

\section{Discussion}
The verification and scaling experiments show that the numerical pipeline is
functionally correct and produces reproducible artifacts. The Poisson refinement
study confirms expected error reduction, while repeated strong-scaling runs
illustrate a practical compute-communication trade-off: gains are clear up to
4 ranks for this workload, but the 8-rank regime becomes variance-heavy and
overhead-dominated on the current machine.

The visualization workflow is also now reliable: each solver writes ParaView-ready
\texttt{.pvd/.vtp} data, and report figures are generated from those exports rather
than hand-built plotting paths. This directly satisfies the original project
constraint that simulation outputs must be useful for ParaView-driven interpretation.

Limitations include short wall-clock solves for some test cases, which can amplify
runtime noise, and the use of representative snapshots for heat/advection figures
instead of full animation panels.

\section{Conclusion and Future Work}
This project now has an end-to-end PETSc-to-ParaView pipeline with validated
numerical behavior, usable performance characterization, and report-ready visuals.
The workflow is immediately portable to Overleaf by keeping figures in a single
directory and including them by filename.

Natural next steps are: (i) larger scaling workloads and additional repetitions
for tighter confidence intervals, (ii) richer time-series visualization products
(animation strips), and (iii) extension to true 3D discretizations and more
advanced preconditioners.

\end{document}
